\paragraph{Block Architecture} -:
The block architecture is like:
\begin{verbatim}
Wishbone
   |
Register File + IRQ Logic (SRB)
   |
QSPI Kernel FSM
   |
TX FIFO / RX FIFO
   |
SPI PHY (SCK / CS / IO[3:0])
\end{verbatim}

The kernel is fully state-driven and does not depend on external events such as interrupts or bus timing.

Writing bit 0 of \asqspiCtrlRegExt\textbf{.START} initiates a transaction if not busy.

\paragraph{Kernel State Machine} -:
The kernel operates in the following states:

\begin{itemize}
  \item IDLE
  \item CMD\_PHASE
  \item ADDR\_PHASE
  \item DUMMY\_PHASE
  \item DATA\_PHASE
  \item DONE
\end{itemize}

State transitions occur only on the internal kernel clock.

\paragraph{Start Semantics:} -:
A transaction is started by:

\[
start\_s = (QSPI\_CTRL.START == 1) \land (BUSY == 0)
\]

The START bit is self-clearing.

\paragraph{START Write while BUSY} -:
Writing the START bit while BUSY is asserted is a no-operation (NOP):
the write is silently discarded and no error is raised. This simplifies
driver implementation (no mandatory pre-check required), but SW is responsible
for polling BUSY or waiting for the DONE interrupt before initiating a
new transaction if ordering is required.

The formal start condition remains:
\[
  start\_s = (\asqspiCtrlRegExt\textbf{.START} == 1) \land (\asqspiStatusRegExt\textbf{.BUSY} == 0)
\]


\paragraph{Error Conditions} -:
Errors are latched in \asqspiStatusRegExt\textbf{.ERROR} and \asqspiRISRegExt\textbf{.ERROR}:
\begin{itemize}
  \item Illegal command
  \item FIFO overflow/underflow
  \item XIP conflict
  \item Timeout
\end{itemize}
Errors persist until cleared via \asqspiICRRegExt.

\paragraph{Execute-In-Place (XIP) Mode} -:
When the XIP bit in \asqspiCtrlRegExt\ is set, the QSPI controller automatically
services read requests from the CPU bus without explicit SW intervention.

\begin{itemize}
  \item XIP is only supported for \emph{read} accesses. Write accesses from the
    bus while XIP is active are not forwarded to the flash and cause an ERROR.
  \item The opcode for XIP reads is taken from \asqspiCmdRegExt\textbf{.OPCODE}.
    SW must configure this register with a suitable read command before
    enabling XIP (e.g.\ 0x6B for Quad Output Fast Read).
  \item The number of dummy cycles is taken from \asqspiDummyRegExt\textbf{.DUMMY\_CYC}.
  \item The ADDR\_LEN, QUAD, DUAL, CPOL, CPHA, and CLKDIV settings apply as
    configured; these must be stable before XIP is enabled.
  \item While XIP is active, the START bit is ignored (writing START has no effect).
  \item XIP is terminated by clearing the XIP bit. The current flash transaction,
    if any, is completed before the controller returns to IDLE.
  \item An XIP conflict error is raised when SW attempts to write
    \asqspiCtrlRegExt\textbf{.START} while XIP is active. This sets the
    ERROR bit in \asqspiStatusRegExt\ and in \asqspiRISRegExt.
\end{itemize}

Typical SW sequence to enable XIP:
\begin{enumerate}
  \item Configure CLKDIV, CPOL, CPHA, QUAD/DUAL, ADDR\_LEN as needed.
  \item Write the desired read opcode to \asqspiCmdRegExt.
  \item Write the required dummy cycles to \asqspiDummyRegExt.
  \item Set XIP bit in \asqspiCtrlRegExt. XIP becomes active immediately.
\end{enumerate}

\paragraph{Design Rules} -:
\begin{itemize}
  \item The kernel FSM is fully state-driven.
  \item No FSM state depends on interrupt signals.
  \item Interrupts never control hardware behavior.
  \item All external behavior is driven by register state only.
\end{itemize}


Figure \ref{fig:asqspi02} shows the state machine of the QSPI.